\documentclass[12pt,a4paper]{article}
\usepackage[UTF8]{ctex}
\usepackage{amsmath,amssymb,amsthm}
\usepackage{geometry}
\usepackage{bm}

\geometry{margin=2.5cm}

\title{旋转坐标系中向量时间导数的推导：$\omega_b \times \bm{L}_b$项的来源}
\author{第8章补充说明}
\date{}

\begin{document}

\maketitle

\section{问题提出}

\textbf{问题}：为什么在旋转坐标系中，向量$\bm{v}$在惯性系中的时间导数包含$\omega \times \bm{v}$这一项？

\textbf{具体问题}：在欧拉方程的推导中，为什么有：
\begin{equation}
\left.\frac{d\bm{L}_b}{dt}\right|_{\text{inertial}} = \left.\frac{d\bm{L}_b}{dt}\right|_{\text{body}} + \omega_b \times \bm{L}_b
\end{equation}

其中$\omega_b \times \bm{L}_b$这一项是怎么来的？

\section{基本公式}

\subsection{旋转坐标系中的时间导数公式}

对于在旋转坐标系中表示的向量$\bm{v}$，其在惯性系中的时间导数为：
\begin{equation}
\left.\frac{d\bm{v}}{dt}\right|_{\text{inertial}} = \left.\frac{d\bm{v}}{dt}\right|_{\text{body}} + \omega \times \bm{v}
\label{eq:rotating_derivative}
\end{equation}

其中：
\begin{itemize}
\item $\left.\frac{d\bm{v}}{dt}\right|_{\text{inertial}}$：在惯性系中观察到的向量变化率
\item $\left.\frac{d\bm{v}}{dt}\right|_{\text{body}}$：在旋转坐标系中观察到的向量变化率
\item $\omega$：旋转坐标系相对于惯性系的角速度
\item $\omega \times \bm{v}$：由于坐标系旋转产生的额外项（科里奥利项）
\end{itemize}

\section{几何直观}

\subsection{简单例子：固定向量}

考虑一个固定在旋转坐标系中的向量$\bm{v}$（在体坐标系中不变）。

\textbf{情况}：假设旋转坐标系绕$z$轴以角速度$\omega$旋转，向量$\bm{v}$固定在坐标系中。

\textbf{观察}：
\begin{itemize}
\item 在体坐标系中：$\bm{v}$不变，$\left.\frac{d\bm{v}}{dt}\right|_{\text{body}} = 0$
\item 在惯性系中：由于坐标系旋转，$\bm{v}$的方向在改变
\end{itemize}

\textbf{几何分析}：

在时间$\Delta t$内，坐标系旋转了$\omega \Delta t$的角度。向量$\bm{v}$在惯性系中的变化为：
\begin{equation}
\Delta \bm{v} = \bm{v}(t + \Delta t) - \bm{v}(t)
\end{equation}

由于旋转，$\bm{v}$的变化方向垂直于$\omega$和$\bm{v}$所在的平面，大小为：
\begin{equation}
|\Delta \bm{v}| \approx |\bm{v}| \sin\theta \cdot \omega \Delta t
\end{equation}

其中$\theta$是$\omega$和$\bm{v}$之间的夹角。

\textbf{方向}：$\Delta \bm{v}$的方向由$\omega \times \bm{v}$给出（右手定则）。

\textbf{因此}：
\begin{equation}
\left.\frac{d\bm{v}}{dt}\right|_{\text{inertial}} = \lim_{\Delta t \to 0} \frac{\Delta \bm{v}}{\Delta t} = \omega \times \bm{v}
\end{equation}

\subsection{一般情况：变化的向量}

如果向量$\bm{v}$不仅在坐标系中变化，坐标系本身也在旋转，那么总的变化率是两项之和：
\begin{enumerate}
\item 在体坐标系中观察到的变化：$\left.\frac{d\bm{v}}{dt}\right|_{\text{body}}$
\item 由于坐标系旋转产生的变化：$\omega \times \bm{v}$
\end{enumerate}

因此：
\begin{equation}
\left.\frac{d\bm{v}}{dt}\right|_{\text{inertial}} = \left.\frac{d\bm{v}}{dt}\right|_{\text{body}} + \omega \times \bm{v}
\end{equation}

\section{严格数学推导}

\subsection{坐标系设置}

设：
\begin{itemize}
\item $\{I\}$：惯性坐标系（固定）
\item $\{B\}$：旋转坐标系（体坐标系）
\item $R(t)$：从$\{B\}$到$\{I\}$的旋转矩阵（时间相关）
\item $\omega$：$\{B\}$相对于$\{I\}$的角速度（在$\{B\}$中表示）
\end{itemize}

\subsection{向量在不同坐标系中的表示}

向量$\bm{v}$在$\{B\}$中的表示为$\bm{v}_b$，在$\{I\}$中的表示为$\bm{v}_i$：
\begin{equation}
\bm{v}_i = R(t) \bm{v}_b
\label{eq:vector_transform}
\end{equation}

\subsection{对时间求导}

对式(\ref{eq:vector_transform})两边对时间求导：
\begin{align}
\frac{d\bm{v}_i}{dt} &= \frac{d}{dt}[R(t) \bm{v}_b] \\
&= \dot{R}(t) \bm{v}_b + R(t) \dot{\bm{v}}_b
\label{eq:derivative_expanded}
\end{align}

\subsection{旋转矩阵的导数}

旋转矩阵的导数满足：
\begin{equation}
\dot{R}(t) = R(t) [\omega]_{\times}
\label{eq:rotation_derivative}
\end{equation}

其中$[\omega]_{\times}$是$\omega$的反对称矩阵：
\begin{equation}
[\omega]_{\times} = \begin{bmatrix}
0 & -\omega_z & \omega_y \\
\omega_z & 0 & -\omega_x \\
-\omega_y & \omega_x & 0
\end{bmatrix}
\end{equation}

\textbf{性质}：$[\omega]_{\times} \bm{v} = \omega \times \bm{v}$

\subsection{代入并简化}

将式(\ref{eq:rotation_derivative})代入式(\ref{eq:derivative_expanded})：
\begin{align}
\frac{d\bm{v}_i}{dt} &= R(t) [\omega]_{\times} \bm{v}_b + R(t) \dot{\bm{v}}_b \\
&= R(t) ([\omega]_{\times} \bm{v}_b + \dot{\bm{v}}_b) \\
&= R(t) (\omega \times \bm{v}_b + \dot{\bm{v}}_b)
\end{align}

\textbf{关键观察}：这个结果是在惯性系$\{I\}$中的表示。要得到在体坐标系$\{B\}$中的表示，需要左乘$R^T$：
\begin{align}
R^T \frac{d\bm{v}_i}{dt} &= R^T R(t) (\omega \times \bm{v}_b + \dot{\bm{v}}_b) \\
&= \omega \times \bm{v}_b + \dot{\bm{v}}_b
\end{align}

\textbf{因此}，在体坐标系$\{B\}$中，向量在惯性系中的时间导数表示为：
\begin{equation}
\left.\frac{d\bm{v}}{dt}\right|_{\text{inertial}} = \dot{\bm{v}}_b + \omega \times \bm{v}_b
\end{equation}

其中：
\begin{itemize}
\item $\dot{\bm{v}}_b = \left.\frac{d\bm{v}}{dt}\right|_{\text{body}}$：在体坐标系中观察到的变化率
\item $\omega \times \bm{v}_b$：由于坐标系旋转产生的额外项
\end{itemize}

\textbf{证毕}。

\section{另一种推导方法：使用旋转矩阵的指数表示}

\subsection{旋转矩阵的指数表示}

旋转矩阵可以表示为：
\begin{equation}
R(t) = \exp([\theta(t)]_{\times})
\end{equation}

其中$\theta(t)$是旋转角度向量，满足$\dot{\theta}(t) = \omega(t)$。

\subsection{对时间求导}

\begin{align}
\dot{R}(t) &= \frac{d}{dt}[\exp([\theta(t)]_{\times})] \\
&= \exp([\theta(t)]_{\times}) [\dot{\theta}(t)]_{\times} \\
&= R(t) [\omega(t)]_{\times}
\end{align}

这与前面的结果一致。

\subsection{应用到向量导数}

\begin{align}
\frac{d}{dt}[R(t) \bm{v}_b] &= \dot{R}(t) \bm{v}_b + R(t) \dot{\bm{v}}_b \\
&= R(t) [\omega]_{\times} \bm{v}_b + R(t) \dot{\bm{v}}_b \\
&= R(t) (\omega \times \bm{v}_b + \dot{\bm{v}}_b)
\end{align}

在体坐标系中表示：
\begin{equation}
R^T \frac{d}{dt}[R(t) \bm{v}_b] = \omega \times \bm{v}_b + \dot{\bm{v}}_b
\end{equation}

\section{应用到角动量}

\subsection{角动量在体坐标系中的表示}

角动量$\bm{L}_b = I_b \omega_b$在体坐标系$\{b\}$中表示。

\subsection{在惯性系中的时间导数}

根据旋转坐标系中的时间导数公式：
\begin{equation}
\left.\frac{d\bm{L}_b}{dt}\right|_{\text{inertial}} = \left.\frac{d\bm{L}_b}{dt}\right|_{\text{body}} + \omega_b \times \bm{L}_b
\end{equation}

\subsection{展开}

\begin{align}
\left.\frac{d\bm{L}_b}{dt}\right|_{\text{inertial}} &= \frac{d}{dt}(I_b \omega_b) + \omega_b \times (I_b \omega_b) \\
&= I_b \dot{\omega}_b + \frac{dI_b}{dt} \omega_b + \omega_b \times (I_b \omega_b)
\end{align}

由于$I_b$在体坐标系中是常数（刚体的形状和质量分布不变），$\frac{dI_b}{dt} = 0$。

\textbf{因此}：
\begin{equation}
\left.\frac{d\bm{L}_b}{dt}\right|_{\text{inertial}} = I_b \dot{\omega}_b + \omega_b \times (I_b \omega_b)
\end{equation}

\subsection{角动量定理}

根据角动量定理：
\begin{equation}
m_b = \left.\frac{d\bm{L}_b}{dt}\right|_{\text{inertial}}
\end{equation}

\textbf{因此}：
\begin{equation}
m_b = I_b \dot{\omega}_b + \omega_b \times (I_b \omega_b)
\end{equation}

这就是欧拉方程。

\section{物理意义}

\subsection{为什么需要$\omega \times \bm{v}$项？}

\textbf{原因}：即使向量在体坐标系中不变，由于坐标系本身在旋转，向量在惯性系中的方向也在改变。

\textbf{例子}：
\begin{itemize}
\item 考虑一个固定在旋转圆盘上的向量
\item 在圆盘坐标系中，向量不变
\item 但在惯性系中，向量的方向在旋转
\item 这个方向变化由$\omega \times \bm{v}$描述
\end{itemize}

\subsection{科里奥利效应}

$\omega \times \bm{v}$项是科里奥利效应的体现：
\begin{itemize}
\item 当物体在旋转坐标系中运动时，会受到额外的"虚拟力"
\item 这个力垂直于运动方向和旋转轴
\item 在动力学方程中表现为$\omega \times \bm{v}$项
\end{itemize}

\subsection{在角动量中的应用}

对于角动量$\bm{L}_b = I_b \omega_b$：
\begin{itemize}
\item 即使$\bm{L}_b$在体坐标系中不变（$\dot{\bm{L}}_b = 0$），由于坐标系旋转，$\bm{L}_b$在惯性系中的方向在改变
\item 这个方向变化由$\omega_b \times \bm{L}_b$描述
\item 这就是为什么需要额外的力矩来维持旋转
\end{itemize}

\section{数值验证}

\subsection{简单例子}

考虑一个绕$z$轴以角速度$\omega = \begin{bmatrix} 0 \\ 0 \\ \omega_z \end{bmatrix}$旋转的坐标系。

固定向量$\bm{v} = \begin{bmatrix} v_x \\ 0 \\ 0 \end{bmatrix}$（在体坐标系中沿$x$轴）。

\textbf{在体坐标系中}：$\dot{\bm{v}}_b = 0$

\textbf{在惯性系中}：
\begin{align}
\omega \times \bm{v} &= \begin{bmatrix} 0 \\ 0 \\ \omega_z \end{bmatrix} \times \begin{bmatrix} v_x \\ 0 \\ 0 \end{bmatrix} \\
&= \begin{bmatrix} 0 \\ \omega_z v_x \\ 0 \end{bmatrix}
\end{align}

\textbf{物理意义}：向量在惯性系中绕$z$轴旋转，$y$方向有速度分量。

\section{总结}

\subsection{主要结论}

\begin{enumerate}
\item \textbf{基本公式}：在旋转坐标系中，向量在惯性系中的时间导数为：
\begin{equation}
\left.\frac{d\bm{v}}{dt}\right|_{\text{inertial}} = \left.\frac{d\bm{v}}{dt}\right|_{\text{body}} + \omega \times \bm{v}
\end{equation}

\item \textbf{$\omega \times \bm{v}$项的来源}：
\begin{itemize}
\item 即使向量在体坐标系中不变，由于坐标系旋转，向量在惯性系中的方向在改变
\item 这个方向变化由$\omega \times \bm{v}$描述
\item 这是旋转坐标系效应的直接结果
\end{itemize}

\item \textbf{应用到角动量}：
\begin{equation}
\left.\frac{d\bm{L}_b}{dt}\right|_{\text{inertial}} = I_b \dot{\omega}_b + \omega_b \times (I_b \omega_b)
\end{equation}

\item \textbf{物理意义}：$\omega_b \times \bm{L}_b$项表示由于坐标系旋转产生的角动量方向变化，这是陀螺效应的体现。
\end{enumerate}

\subsection{关键要点}

\begin{itemize}
\item 旋转坐标系中的时间导数必须考虑坐标系旋转的影响
\item $\omega \times \bm{v}$项是几何和物理的必然结果
\item 这个项在动力学方程中表现为科里奥利效应
\item 在角动量定理中，这个项导致需要额外的力矩来维持旋转
\end{itemize}

\subsection{应用场景}

\begin{itemize}
\item 刚体动力学（欧拉方程）
\item 机器人动力学（递归牛顿-欧拉算法）
\item 航天器姿态控制
\item 陀螺仪和导航系统
\item 旋转机械系统分析
\end{itemize}

\end{document}

